% ===== PRÉAMBULE COMMUN — À COPIER TEL QUEL EN TÊTE DE CHAQUE FICHIER .tex =====
% Ne modifier QUE :
%   - les deux lignes \fancyhead[L] et \fancyhead[R]
%   - le bloc titre \begin{center}...\end{center} juste après \begin{document}
\documentclass[11pt,a4paper]{article}
\usepackage[utf8]{inputenc}
\usepackage[T1]{fontenc}
\usepackage{lmodern}
\usepackage[french]{babel}
\usepackage{amsmath,amssymb,mathtools}
\usepackage{icomma}
\usepackage{xcolor}
\usepackage{tikz}
\usepackage{pgfplots}
\usepackage[most]{tcolorbox}
\usepackage{enumitem}
\usepackage{array,booktabs,colortbl}
\usepackage[a4paper,margin=2.2cm,headheight=26pt,headsep=10pt]{geometry}
\usepackage{fancyhdr}
\usepackage[hidelinks]{hyperref}
\usetikzlibrary{arrows.meta,calc,angles,quotes,positioning,decorations.markings,intersections,backgrounds}
\pgfplotsset{compat=1.18}

\definecolor{primary}{RGB}{226,74,88}
\definecolor{primarydark}{RGB}{150,28,42}
\definecolor{methcol}{RGB}{40,90,150}
\definecolor{excol}{RGB}{0,135,90}
\definecolor{defcol}{RGB}{0,114,178}
\definecolor{attncol}{RGB}{200,40,55}
\definecolor{thmcol}{RGB}{0,140,120}
\definecolor{courbe}{RGB}{32,110,200}
\definecolor{courbeder}{RGB}{210,70,40}
\definecolor{tang}{RGB}{0,150,90}
\definecolor{grille}{RGB}{200,205,215}
\definecolor{plus}{RGB}{200,60,60}
\definecolor{moins}{RGB}{30,90,200}

\tcbset{boxbase/.style={enhanced,breakable,boxrule=0.9pt,arc=2.5pt,
  left=3mm,right=3mm,top=2.3mm,bottom=2.3mm,fonttitle=\bfseries,coltitle=white}}
\newtcolorbox{cle}[1][]{boxbase,colback=defcol!6,colframe=defcol,
  title={\textsc{À retenir}\ifx\relax#1\relax\relax\else\textmd{ : \itshape #1}\fi}}
\newtcolorbox{methode}[1][]{boxbase,colback=methcol!7,colframe=methcol,sharp corners,
  title={\textsc{Méthode}\ifx\relax#1\relax\relax\else\textmd{ --- \itshape #1}\fi}}
\newtcolorbox{exemplebox}[1][]{boxbase,colback=excol!5,colframe=excol,
  title={\textsc{Exemple}\ifx\relax#1\relax\relax\else\textmd{ : \itshape #1}\fi}}
\newtcolorbox{piege}[1][]{boxbase,colback=attncol!6,colframe=attncol,
  title={\textsc{Attention}\ifx\relax#1\relax\relax\else\textmd{ : \itshape #1}\fi}}
\newtcolorbox{exo}[1][]{boxbase,colback=primary!4,colframe=primary,
  title={\textsc{Exercice}\ifx\relax#1\relax\relax\else\textmd{~#1}\fi}}
\newtcolorbox{corr}[1][]{boxbase,colback=thmcol!5,colframe=thmcol,
  title={\textsc{Corrigé}\ifx\relax#1\relax\relax\else\textmd{~#1}\fi}}

\newcommand{\R}{\mathbb{R}}
\newcommand{\N}{\mathbb{N}}
\newcommand{\ds}{\displaystyle}
\newcommand{\tg}{\operatorname{tg}}
\newcommand{\cotg}{\operatorname{cotg}}
\newcommand{\arctg}{\operatorname{arctg}}
\newcommand{\arccotg}{\operatorname{arccotg}}
\newcommand{\limh}{\lim_{h\to 0}}

\pagestyle{fancy}\fancyhf{}
\fancyhead[L]{\small\textcolor{primarydark}{\textbf{Thème 3} — Fonctions composées}}
\fancyhead[R]{\small\textcolor{primarydark}{\textsc{Corrigé — Difficile}}}
\fancyfoot[C]{\small\thepage}
\renewcommand{\headrulewidth}{0.8pt}

\begin{document}

\begin{center}
{\LARGE\bfseries\color{primarydark} Plusieurs couches \& fonctions inconnues — corrigé}\\[2pt]
{\color{primary}\rule{0.5\textwidth}{1.2pt}}\\[4pt]
{\small\itshape Niveau difficile}
\end{center}
\vspace{2mm}

\begin{corr}[1 --- $y=\sin^2(\sqrt x+1)$]
\textbf{a)} Externe = le carré. Avec $(F^2)'=2F\cdot F'$ appliqué à $F=\sin(\sqrt x+1)$ :
\[
y'=2\sin(\sqrt x+1)\cdot\big(\sin(\sqrt x+1)\big)'.
\]
\textbf{b)} Pour $\sin(\sqrt x+1)$ : couche externe $\sin$, intérieur $U=\sqrt x+1$ avec
$U'=\dfrac{1}{2\sqrt x}$. Donc $(\sin U)'=U'\cos U$ :
\[
\big(\sin(\sqrt x+1)\big)'=\frac{1}{2\sqrt x}\cos(\sqrt x+1).
\]
\textbf{c)} On assemble :
\[
y'=2\sin(\sqrt x+1)\cdot\frac{1}{2\sqrt x}\cos(\sqrt x+1)
 =\frac{1}{2\sqrt x}\cdot\underbrace{2\sin(\sqrt x+1)\cos(\sqrt x+1)}_{=\,\sin(2(\sqrt x+1))}.
\]
D'où, avec $2\sin\theta\cos\theta=\sin 2\theta$ ($\theta=\sqrt x+1$) :
\[
\boxed{\,y'=\frac{\sin\big(2(\sqrt x+1)\big)}{2\sqrt x}\,}\qquad\text{pour } x>0
\]
(on exclut $x=0$ car $\sqrt x$ y est non dérivable et le dénominateur s'annule).
\end{corr}

\begin{corr}[2 --- $y=\sqrt{\ln(\cos x+2)}$]
\textbf{a)} On écrit $y=\big(\ln(\cos x+2)\big)^{1/2}$. Couche puissance $\tfrac12$, avec
$F=\ln(\cos x+2)$ et $(F^{1/2})'=\dfrac{F'}{2\sqrt F}$ :
\[
y'=\frac{\big(\ln(\cos x+2)\big)'}{2\sqrt{\ln(\cos x+2)}}.
\]
\textbf{b)} Couche $\ln$ : $U=\cos x+2$, $U'=-\sin x$, donc $(\ln U)'=\dfrac{U'}{U}$ :
\[
\big(\ln(\cos x+2)\big)'=\frac{-\sin x}{\cos x+2}.
\]
\textbf{c)} On assemble :
\[
y'=\frac{1}{2\sqrt{\ln(\cos x+2)}}\cdot\frac{-\sin x}{\cos x+2}
 =\boxed{\,\frac{-\sin x}{2(\cos x+2)\sqrt{\ln(\cos x+2)}}\,}.
\]
(Comme $\cos x+2\ge 1$, on a $\ln(\cos x+2)\ge 0$ ; il faut $\ln(\cos x+2)>0$,
soit $\cos x+2>1$, i.e. $\cos x\neq -1$.)
\end{corr}

\begin{corr}[3 --- fonction inconnue $f$]
\textbf{a)} $f^2=f\cdot f$. Règle du produit : $(f\cdot f)'=f'f+f f'=2f'f$. (Ou bien
$(U^2)'=2UU'$ avec $U=f$.) Donc $\big(f^2\big)'=2f'f$.\\[2pt]
\textbf{b)} $\big(\ln|f|\big)'=\dfrac{f'}{f}$ (formule $(\ln|U|)'=U'/U$). En multipliant par la
constante $3$ : $\big(3\ln|f|\big)'=\dfrac{3f'}{f}$.\\[2pt]
\textbf{c)} $f(x)=x^2+1$, donc $f'(x)=2x$.
\[
\big(f^2\big)'=2f'f=2(2x)(x^2+1)=4x(x^2+1),\qquad
\big(3\ln|f|\big)'=\frac{3f'}{f}=\frac{6x}{x^2+1}.
\]
\end{corr}

\begin{corr}[4 --- fonction inconnue $u$]
\textbf{a)} $v(x)=u(3x)$ : externe $u$, intérieur $U=3x$ avec $U'=3$. La chaîne donne
$v'(x)=u'(3x)\cdot 3=3u'(3x)$.\\
Si $u(x)=x^2$, alors $u'(x)=2x$, donc $u'(3x)=2(3x)=6x$ et $v'(x)=3\cdot 6x=18x$.\\
Vérification directe : $v(x)=u(3x)=(3x)^2=9x^2$, d'où $v'(x)=18x$. \checkmark\\[2pt]
\textbf{b)} $v(x)=u\big(\ln(x+1)\big)$ : externe $u$, intérieur $U=\ln(x+1)$ avec
$U'=\dfrac{1}{x+1}$. Donc
\[
v'(x)=u'\big(\ln(x+1)\big)\cdot\frac{1}{x+1}=\frac{u'\big(\ln(x+1)\big)}{x+1}\qquad(x>-1).
\]
\end{corr}

\end{document}
